\subsubsection*{推论 6.7.}
在 6.6 的假设下, 假设 $X$ 满足条件 (LV). 则:

\medskip
\noindent
\textup{6.7.1.} $\pi_1(D-D\cap\check X)$ 在 $E^{n-1}(X_{\bar\eta},\mathbb Q_\ell)$ 上的表示绝对不可约.

\medskip
\noindent
\textup{6.7.2.} 若 $n$ 为偶数, 则 $\pi_1(D-D\cap\check X)$ 的每个有限指数子群在 $E^{n-1}(X_{\bar\eta},\mathbb Q_\ell)$ 上的作用均绝对不可约.

\begin{editorialnote}
这里的证明表明消灭上同调没有非零真不变子空间. 本讲的假设也允许消灭上同调为零 (例如 5.9 的二次曲面); 按不可约表示须非零的通常约定, 6.7.1、6.7.2 的“绝对不可约”结论须在该空间非零时理解. 零空间的情形则只有上述关于不变子空间的空泛断言. 此处明确这一边界情形, 不另加源文未有的几何假设.
\end{editorialnote}

\subsubsection*{证明.}
可以将所论模换成其扭转
\[
E=E^{n-1}(X_{\bar\eta},\mathbb Q_\ell(m)).
\]
设 $\bar E=E\otimes_{\mathbb Q_\ell}\bar{\mathbb Q}_\ell$, 并选取 $0\ne\chi\in\bar E$. 为证明 6.7.1, 须说明 $\bar E$ 中由 $\chi$ 及其在 $\pi$ 下的共轭元生成的 $\bar{\mathbb Q}_\ell$-线性子空间就是整个 $\bar E$. 由 6.6.1 和 6.6.2, 只须证明此子空间包含某个消灭圈 $\delta_i$.

由于 $X$ 满足 (LV), 杯积型在下式上的限制
\[
H^{n-1}(X,\mathbb Q_\ell(m))
\lhook\joinrel\longrightarrow
H^{n-1}(X_{\bar\eta},\mathbb Q_\ell(m))
\]
非退化; 因而在正交补 $E^{n-1}(X_{\bar\eta},\mathbb Q_\ell(m))$ 上的限制也非退化. 由 6.6.1, 因此至少存在一个 $i$, 使得
\[
\chi\cdot\delta_i\ne0.
\]
对于这个 $i$, 6.5 给出 $s_i$ 处惯性群的一个元素 $\sigma_i$, 使得
\begin{equation}\tag{6.7.3}
\sigma_i(\chi)-\chi
=\text{下式的一个非零倍数:}\delta_i.
\end{equation}
这就证明了 6.7.1.

对于 6.7.2, $n-1$ 为奇数, 所以 $\delta_i\cdot\delta_i=0$. 迭代使用 6.5.2, 对每个 $k$, 得
\begin{equation}\tag{6.7.4}
\sigma_i^k(\chi)
=\chi-k(-1)^m\sigma_i(\chi\cdot\delta_i)\delta_i.
\end{equation}
\begin{editorialnote}
原印本及英译底本在 (6.7.4) 右边写作 $\chi+k\sigma_i(\chi\cdot\delta_i)\delta_i$, 遗失了 (6.5.2.1) 中的因子 $-(-1)^m$. 因 $\delta_i\cdot\delta_i=0$, 令 $N=\sigma_i-1$ 则 $N^2=0$, 从而 $\sigma_i^k\chi=\chi+kN\chi$; 此处据此补回该因子. 后面的不变子空间论证只用此系数非零, 不受此订正影响.
\end{editorialnote}

因而 $E$ 的 $\mathbb Q_\ell$-子空间 $M$ 在 $\pi$ 的作用下稳定, 当且仅当它在某个有限指数子群的作用下稳定: 两个条件均等价于
\[
\chi\in M
\quad\Longrightarrow\quad
\delta_i\in M\quad\text{对每个满足下式的 $i$: }
\chi\cdot\delta_i\ne0.
\]
因此断言归结为 6.7.1.

\subsection*{参考文献}

\begin{description}
\item[{[1]}] P. Deligne, \emph{Théorème de Lefschetz et critères de
dégénérescence de suites spectrales}, Publ. Math. IHÉS \textbf{35}
(1968), 259--278.

\item[{[2]}] S. Kleiman, \emph{Algebraic cycles and the Weil
conjectures}, in \emph{Dix exposés sur la cohomologie des schémas},
North-Holland, 1968.

\item[{[3]}] S. Lefschetz, \emph{L'analyse situs et la géométrie
algébrique}, Gauthier-Villars, Paris, 1924.

\item[{[4]}] D. Lieberman, \emph{Higher Picard Varieties}, Amer. J.
Math. \textbf{90} (4) (1968), 1165--1199.\footnote{印本将期刊卷号写成 15; 该文发表于第 90 卷.}

\item[{[5]}] D. Mumford, 未刊.

\item[{[6]}] A. Weil, \emph{Variétés kählériennes}, Hermann, Paris,
1958.
\end{description}

